\lstdefinelanguage{json}{
    basicstyle=\normalfont\ttfamily,
    numbers=left,
    numberstyle=\scriptsize,
    stepnumber=1,
    numbersep=8pt,
    showstringspaces=false,
    breaklines=true,
    frame=lines,
    backgroundcolor=\color{gray!10},
    string=[s]{"}{"},
    comment=[l]{:\ "},
    morecomment=[l]{:"},
    literate=
     *{0}{{{\color{nblue}0}}}{1}
      {1}{{{\color{nblue}1}}}{1}
      {2}{{{\color{nblue}2}}}{1}
      {3}{{{\color{nblue}3}}}{1}
      {4}{{{\color{nblue}4}}}{1}
      {5}{{{\color{nblue}5}}}{1}
      {6}{{{\color{nblue}6}}}{1}
      {7}{{{\color{nblue}7}}}{1}
      {8}{{{\color{nblue}8}}}{1}
      {9}{{{\color{nblue}9}}}{1},
}


\chapter{Estudio previo}\label{cap:planificación}

\section{Introducción}
Breve introducción al capítulo

\subsection{El estándar ERC-721}
Según la documentación oficial \cite{openzeppelin_erc721}, el ERC-721 es el estándar para representar 
la propiedad de tokens no fungibles (NFTs), donde cada token es único e indivisible. En nuestra arquitectura, 
el NFT actúa como la clave primaria y el certificado de propiedad irrefutable de cada carta coleccionable.

\subsubsection{Estructura de Datos}
A diferencia de los SGBD tradicionales que usan tablas relacionales y claves foráneas, el ERC-721 asocia 
la propiedad mediante diccionarios de datos nativos de la blockchain (\verb|mapping|).

A través de funciones estándar como \verb|ownerOf(tokenId)|, el contrato inteligente actúa como nuestro motor 
de base de datos, permitiéndonos consultar qué dirección criptográfica (usuario) posee un ID de carta 
específico de forma instantánea y segura.

\subsubsection{Estrategia de Almacenamiento y Costes}
Almacenar grandes cantidades de datos en la blockchain, como imágenes de alta calidad, tiene un coste de 
transacción alto y prohibitivo. La propia documentación de OpenZeppelin advierte sobre esto y sugiere el uso 
de sistemas como IPFS.

Por ello, aplicamos una división de datos:
\begin{itemize}
    \item \textbf{Datos On-Chain}. \\ Solo se persisten en la blockchain los datos considerados críticos que garantizan la titularidad:
    \begin{itemize}
        \item \textbf{tokenId:} ID de la carta. % Corregido: decía tokenIf
        \item \textbf{owner:} Dirección del propietario.
        \item \textbf{tokenURI:} Un puntero o enlace al alojamiento externo.
    \end{itemize}

    \item \textbf{Datos Off-Chain}. \\ El \texttt{tokenURI} se diseña para que resuelva a un documento JSON alojado en la 
    red descentralizada IPFS. Este archivo contiene los datos pesados de la carta: % Corregido: comillas Markdown eliminadas

    \begin{lstlisting}[language=json, caption=Ejemplo de tokenURI apuntando a IPFS]
{
    "name": "Charizard",
    "description": "Carta holografica primera edicion",
    "image": "ipfs://<CID-de-la-imagen>"
}
    \end{lstlisting}
\end{itemize}

\section{Objetivos}

\section{Metodología}

\section{Planificación}

\section{Presupuesto}

\section{Conclusiones}